\subsection{Discordant analysis clustering}\label{sec:disc-analysis-cluster}

Discordant zircon populations may contain more than one internally coherent discordant sub-array. When such analyses are evaluated \textit{en masse}, the concordant--discordant comparison (CDC) may be drawn toward a compromise lower-intercept age that reflects the numerically dominant discordant population rather than preserving reproducible subordinate structure. To test for this possibility, discordant analysis clustering is implemented as an optional pre-partitioning step prior to ensemble peak picking.

When clustering is enabled, each discordant analysis is assigned a single sample-level lower-intercept proxy age using concordant-reference anchors in Tera--Wasserburg space. Concordant analyses are first represented by a small set of concordant age populations estimated using one-dimensional Gaussian mixture modelling with Bayesian information criterion selection; in the present implementation, up to six concordant anchor populations are allowed so that samples with more complex concordant structure are not forced into an overly restrictive anchor set. The centroids of these concordant populations are then used as anchor points. For each discordant analysis, the line linking the analysis to each candidate anchor is evaluated across the tested Pb-loss age grid, and the lower-intercept proxy age is taken as the candidate age that minimises the geometric misfit to the corresponding concordia chord. This produces one anchored proxy age per discordant analysis, fixed at the sample level and not reassigned during Monte Carlo resampling.

The resulting discordant proxy ages are clustered once per sample in one dimension using Gaussian mixture modelling, after which the candidate split is filtered by explicit size and separation criteria. A clustered solution is accepted only if at least two groups remain after gating, each retained group contains a minimum proportion of the discordant population, and the retained groups are separated by more than the larger of their robust dispersions. If these criteria are not met, the sample is treated as unclustered. Clustering is therefore used only where the discordant population shows evidence of internally coherent subdivision rather than being imposed as a default interpretation.

If a robust split is accepted, CDC goodness surfaces are evaluated separately for the retained discordant clusters and ensemble peaks are picked from those clustered surfaces using the same support, prominence, width, and merging criteria applied elsewhere in the workflow. Clustered peaks are not reported automatically, however. The unsplit all-discordants CDC surface remains the primary reporting reference: clustering is permitted to refine peak structure only where the corresponding all-discordants surface already preserves reproducible multimodality. If the all-discordants surface is unresolved, monotonic, boundary-dominated, or supports only a single peak, the sample is reported from the unclustered solution and any additional cluster-only peaks are discarded. In this way clustering functions as a constrained pre-partitioning step for structured discordant populations rather than as an independent mechanism for creating additional disturbance ages.

Recovered CDC peaks are treated as candidate disturbance ages whose geological significance must be assessed against petrographic, textural, and regional constraints. For the synthetic benchmarks, clustering should be evaluated using one fixed parameter set across the benchmark suite rather than enabled or disabled on a case-by-case basis, so that the clustered workflow is tested as a reproducible method rather than selected only where it gives a preferred outcome.
